\cotefile{72}
\scanpage{2}
(23 p)

$n$ polyèdres réguliers et représentations \unsure{des|du|de} $G_n$ dans $O(n-1)$ \unsure{.....\,--|......|.....}

(dictionnaire)

% archive sheet number: 1

\scanpage{3}
$n\geqslant1$ fixé, et $n$-hyperpolyèdre combinatoire épinglé universel $ {}_n\Pi$, groupe d'autom.
\[
\unsure{\mathscr G_n|\mathfrak g_n|G_n|\mathfrak G_n}=\{\tau_0,\ldots,\tau_n\mid \tau_i\longleftrightarrow\tau_j\text{ si }0\leq i,j\leq n,\ i,j\text{ non cons.}\}
\]
$S$ schéma de base fixé

\margin{autodual} $C_{pr}$ cat. des réal. géom. régulières \inserted{non dég} projectives de $\Pi$ (autodual)

\margin{autodual} $C'_{pr}$ cat des systèmes $\deleted{(X,\varphi_G:\mathscr G_n\to\operatorname{Aut}_S X)}$ $(X,(\tau_i)_{0\leq i\leq n},\deleted{\unsure{b}}X_*)$ $X$ fibré projectif \inserted{rg $n+1$}, $\tau_i$ $(0\leq i\leq n)$ ps. refl. strictes \inserted{($\Rightarrow h_i,H_i$)}, \deleted{$s_0$ section} $X_*$ drapeau maximal, avec conditions
\[
\begin{aligned}
\text{a)}\ &\tau_i\longleftrightarrow\tau_j\qquad\text{si}\qquad 0\leq i,j\leq n\qquad i,j\text{ non consécutifs.}\\
\text{b)}\ &\tau_i X_j=X_j\quad\text{si }0\leq i\ne j\leq n\\
\text{c)}\ &\tau_i X_i\ne X_i\qquad 0\leq i\deleted{\unsure{,j}}\leq n
\end{aligned}
\]
\margin{\uline{NB} $C_{pr},C'_{pr}$ proviennent par \uline{rigidification} de catégories $\widetilde C_{pr}$ \deleted{\unsure{$C'_{pr}$}} de \uline{représentations} projectives de $\mathscr G_n$ \deleted{sans} plus}
\margin{\uline{NB} Au lieu de se donner $X_*$, il suffit de se donner \uline{un} $X_i$ avec $\tau_jX_i=X_i$ $(j\ne i)$ $\tau_iX_i\ne0$ satisfaisant à conditions ouv. -- p. ex $s_0=X_0$, tel que \deleted{la suite des \unsure{$s$}} \inserted{$(s_0,(h_i))$} engendre \unsure{lin$^{\mathrm{t}}$} $X$}

\margin{non autodual} $C_{aff}$ variante affine de $C_{pr}$

\margin{non autodual} $C'_{aff}$ variante affine de $C'_{pr}$ \inserted{[NB. il n'est plus besoin de se donner $\varphi_G$ \unsure{; idem|; itou} pour $C'_{aff}$].}

(Attention, $\tau_0$ n'est peut-être pas une ps. réflexion affine, mais $\widehat\tau_0$ est une ps. refl. projective (dont le \deleted{hyp} centre peut être exceptionnellement l'hyperplan à l'$\infty$ \ldots\ )
\margin{cas $\beta_0=\mu_0=0$ ! (avéré en car. $2$..$)$}

\margin{autodual} $C_{vect}$ variante vectorielle de $C_{pr}$ (où $X$ \uline{rigidisé} comme $P(E)$, et $s_0$ rigidisé par une section \uline{de} $E$ -- ce qui \deleted{rigidifie} est une donnée de \uline{rigidification} dans une donnée $\in$ Catégorie $\widetilde C_{vect}$, éliminant le groupe $G_m$ d'autom..
\margin{$\widetilde C_{vect},\widetilde C'_{vect}$ sans $G_m$-épinglage}

$C'_{vect}$ \unsure{variante} vectorielle de $C'_{pr}$ (on laisse tomber \unsure{qui est|qu'est} $=$ la \unsure{der. prem.|drapeau}, remplacée par $s_0=a_{-1}\ldots$)

$\deleted{\unsure{\,\sim\,}}C^{\infty}_{\unsure{pr*|pr\ne *}}$ cat des $(C,\varphi_G^C:\mathscr G_n\to\operatorname{Aut} C,b_0)$ où $C$ fibré projectif de rg $n$, $\varphi_G^C$ hom groupe, $h_0$ section de $C$, avec
\margin{non autodual, sauf si $\unsure{\rho|p}\ne\unsure{0|h_0}$ i.e. $c\notin C$ en tte fibre}
\margin{\uline{NB} La donnée de $s_0$ superflue si $\tau_0^C$ stricte i.e. $\mu_0,\beta_0$ non nuls simultanément}
a) les $\tau_i^C$ ps. réflex. $(0\leq i\leq n)$ strictes si $1\leq i\leq n$, associées à $\unsure{h_0|b_0}$ pour $\tau_0^C$ (qui peut être non stricte ($\Rightarrow \deleted{\unsure{\exists}}h_i,H_i$ pour $1\leq i\leq n$ -- mais pas de $h_0$)

% archive sheet number: 2

\scanpage{4}
\[
\begin{aligned}
\text{b)}\ & h_i\quad(0\leq i\leq n)\quad \text{engendrent } C\quad\text{i.e. lin. indép.}\\
\text{c)}\ & h_i\notin H_{i+1}\quad(0\leq i\leq n-1)\quad\text{en tte fibre}
\end{aligned}
\]

\margin{non autodual sauf si $\rho$ inv.}
$\widetilde C^\infty_{vect}$ variante vectorielle $(V,\varphi_{\unsure{G|g|\mathfrak G}}^V,L_0)$, $V$ fibré vect. rg $n+1$, $\varphi_G^V:\unsure{\mathscr G|G|\mathfrak G}\longrightarrow\operatorname{Aut}(V)$, $L_0$ sous-fibré vect. de rg $1$, avec

a) les $\tau_i^V$ $0\leq i\leq n$ ps. réflexions, strictes pour $1\leq i\leq n$, et sub. à $L_0$ si $i=0$ ($\Rightarrow L_i,H_i,L'_i$ si $1\leq i\leq n$)
\[
\begin{aligned}
\text{b)}\ & V=\bigoplus_{0\leq i\leq n} L_i\\
\text{c)}\ & h_i\notin H_{i+1}\qquad(0\leq i\leq n-1)
\end{aligned}
\]
Alors $L_0\simeq L_1^\vee\simeq L_1\simeq L_2^\vee\simeq L_2\ldots\simeq L_n^\vee\simeq L_n$, soit $L$.

On appelle \uline{rigidification} la donnée d'une base de $L$, (p. ex. d'une base $a_0$ de $L_0$, d'où $a_i$ de $L_i$, $a'_i$ de $L'_i$).

$C^\infty_{vect}$ cat. des objets de $\widetilde C^{\unsure{n|\infty}}_{vect}$ rigidifiés

\margin{$n\geqslant 2$, les $\mu_i=-2$}
\margin{non autodual sauf si $\rho$ inv.}
$C^\infty_{pr.\,orth}$ cat des objets $(C,Q,\varphi_G^{C,Q},b_0)$, où $C$, $\deleted{\unsure{\varphi_G^C|\varphi_{\mathfrak G}^C}}$ fibré projectif rg $n$, $Q$ hyper\inserted{$\unsure{\infty}$}quadrique \deleted{la donnée \unsure{revient} par $C^\infty_{prog}$, et $Q$ hyperquadrique inv. de $C$,}
\[
\varphi_G^{C,Q}:\mathscr G_n\longrightarrow\operatorname{Aut}(C,Q),
\]
$b_0$ section de $C$, tels que

a) les $\tau_i^{C,Q}$ $(0\leq i\leq n)$ réfl. orth, strictes pour $1\leq i\leq n$, sub. à $h_0$ si $i=0$ ($\Rightarrow h_i,H_i\quad1\leq i\leq n$)
\[
\begin{aligned}
\text{b)}\ &h_i\quad (0\leq i\leq n)\quad\text{engendrent }C\\
\text{c)}\ &h_i\notin H_{i+1}\quad\text{en tte fibre}\quad(0\leq i\leq n-1)
\end{aligned}
\]

\margin{\uline{NB} \deleted{Dans} le cas \deleted{\unsure{propre}} \inserted{$\beta_0\beta_1$ inv.} (plus \inserted{gén} $\beta_0$ \deleted{\unsure{prem.}} inv., \uline{et} $\deleted{(\beta_1,2)\mathcal O_S=\mathcal O_S}$ \deleted{($\beta_1$ inv. en tte \unsure{pt} de car. $2$)} $(\beta_1,\deleted{\unsure{\beta}}4-\beta_0)\mathcal O_S=\mathcal O_S$ i.e. $\beta_1=\alpha_1+2$ et $4-\beta_0=2-\alpha_0$ non nuls simultanément) b) équivaut au \uline{fait} que sur aucune fibre, il n'y a d'hyperplan \uline{de $C$} invariant par $G$ \ldots\ (cette remarque valant aussi pour $C^\infty_{vect}$, $C^\infty_{pr}$ \ldots. expliciter}
\uline{ReNB} si $\tau_i^C,\tau_{i+1}^C$ réflexions \uline{str.} orth. i.e. $\beta_0\ldots\beta_i$ inv. Alors $h_i\notin H_{i+1}$ en tte fibre équivaut à $\underbrace{\tau_i\tau_{i+1}}_{\rho_i}\ne\tau_{i+1}\tau_i$ \inserted{tte fibre} i.e. $\rho_i^2\ne\mathrm{id}$ en tte fibre. Donc si $\beta$ inv. c) se remplace par ces conditions pour $0\leq i\leq n-1$ \ldots

% archive sheet number: 3

\scanpage{5}
\uline{\unsure{Rere|Re re|Rem} NB} On aura \unsure{néc.|même}$^{\mathrm t}$ $h_0\in\Gamma(C\smallsetminus Q)$ ....

\margin{non autodual sauf si $\beta$ inv.}
$C^\infty_{vect\,orth}$ cat. des syst. $(V,f,\varphi_G^V,a_0)$, où $V$ fibré vectoriel de rg $n+1$, $f$ forme quadratique \deleted{non id. nulle sur aucune fibre}, $\varphi_G^{V\,orth}:\unsure{\mathscr G|G|\mathfrak G}\longrightarrow\operatorname{Aut}(V,f)$, $a_0$ section de $V$, satisfais. à

a) les $\tau_i^V$ réflexions orth \inserted{$(0\leq i\leq n)$}, les $\tau_i$ $(1\leq i\leq n)$ \deleted{\unsure{np|non}} $\ne$ id en chaque fibre ($\Rightarrow L_i,H_i$ pour $1\leq i\leq n$), $\tau_0$ associé à $\mathcal O_S a_0$, $\deleted{(a_0\ne0\text{ en tte}}$ et $\boxed{f(a_0)=1}$ (de sorte que $L_0$ $\mathcal O_S\deleted{\unsure{b_0}}a_0$ ss-fibré vect\ldots)

b), c) comme dessus, avec commentaire correspondant.

On pourrait définir $\widetilde C^\infty_{vect\,orth}$, catégorie des \deleted{systèmes} $(V,f,\varphi_G^V,L_0)$, où \deleted{$(\unsure{V},\unsure{f},\varphi_G^V,L_0)$} \deleted{\unsure{comme} $\varphi_G^V$, $f$} $(V,f)$ fibré vectoriel quadratique \deleted{comme dans $C^\infty_{vect}$, $(V,f)$ fibré vect.} de rg $n+1$, $\varphi_G^V:\mathscr G_n\longrightarrow\operatorname{Aut}(V,f)$, $L_0$ ss-fibré de rg $1$, tels que a) b) c) comme précédemt, mais $f(a_0)=1$ remplacé par $f\mid L_0$ est \uline{\unsure{lisse|libre}}.

Et On trouve $C^\infty_{vect\,orth}$ à partir de $\widetilde C^\infty_{vect\,orth}$ par rigidification par une section $a_0$ de $L_0$ telle que $f(a_0)=1$ (les sections forment un torseur sous $\mu_2$ \ldots\ ).

\diagram{\begin{tikzcd}[ampersand replacement=\&,column sep=small]
C_{pr}\arrow[r,dash,"\approx"]\& C'_{pr}\arrow[r,dash,"\approx"]\& C_{aff}\arrow[r,dash,"\approx"]\& C'_{aff}\arrow[r,dash,"\approx"]\& C_{vect}\arrow[r,dash,"\approx"]\& C'_{vect}\,\Big|\arrow[r,dash,"\approx"]\& C^\infty_{pr}
\end{tikzcd}}
\margin{$n\geqslant2$ (expliciter variante si $n=1$ !)}
\diagram{\begin{tikzcd}[ampersand replacement=\&,column sep=small]
{}\arrow[r,dash,"\approx"]\& C^\infty_{vect}\arrow[r,dash,"\approx"]\& \Big|\,C^\infty_{pr\,orth}\arrow[r,dash,"\approx"]\& C^\infty_{vect\,orth}
\end{tikzcd}}
\margin{si de plus les $\lambda_i=-1$ i.e. les $\tau_i$ des réflexions.}

% archive sheet number: 4

\scanpage{6}
Soit $G_0$ l'image de $\unsure{\mathscr G_n|G_n|\mathfrak G_n}$ dans $\operatorname{Aut}(E,\deleted{\unsure{f},}V)$, ou mieux dans $\operatorname{Aut}(E,f,\pi,\pi')$
\diagram{\begin{tikzcd}[ampersand replacement=\&,column sep=large]
\deleted{G_0}\arrow[r,draw=none,"\deleted{\longrightarrow}" description]\&\operatorname{Aut}(X,C)\deleted{\unsure{Q_X}}\arrow[r]\&\operatorname{Aut}(C,\deleted{\unsure{Q_C}})\\
G_0\arrow[r,hook]\&\operatorname{Aut}(E,V,f)\arrow[u]\arrow[r]\&\operatorname{Aut}(V,\deleted{f_V})\arrow[u]
\end{tikzcd}}
On cherche \deleted{les} $\inserted{3}$ noyaux de
\diagram{\begin{tikzcd}[ampersand replacement=\&,column sep=small,row sep=small]
\&[-1em] \operatorname{Aut}(X,\deleted{\unsure{C}})\arrow[dr]\&\\
G_0\arrow[r]\&\left\{\vphantom{\dfrac{A}{A}}\right.\&\operatorname{Aut}(C)\\
\&\operatorname{Aut}(V)\arrow[ur]\&
\end{tikzcd}}

a) $\operatorname{Ker}(\operatorname{Aut}(E,V,f)\deleted{\unsure{\to}}\longrightarrow\operatorname{Aut}(X))$ est formé des homothéties qui conservent $f$, donc c'est $\mu_2$. En fait, $G_0$ conserve mieux que $(E,f,V)$, il conserve $(E,f,\pi,\pi')$ \inserted{\unsure{pour n'importe $f$|peu importe $f$}}, donc si c'est une homothétie c'est l'identité : $G_0$ opère \uline{fidèlem$^{\mathrm t}$} sur $X$.

b) $\operatorname{Ker}(\operatorname{Aut}(E,f,\pi,\pi')\longrightarrow\operatorname{Aut}(V))$ est formé des morphismes \inserted{$u$} de la forme \deleted{\unsure{$u$}} $\mathrm{id}+\deleted{\unsure{\lambda}}\pi'\otimes a$ $(a\in\Gamma E)$, qui conservent $\pi,\pi',f$. Or $u$ conserve $\pi$ sss $\deleted{\unsure{\lambda}}\rho a=0$ (donc $a=0$ si $\rho$ inv., plus gén. si $\rho$ est \uline{régulier}) (NB $\langle\pi,\pi'\rangle=\rho$), et il conserve $\pi'$ sss :
\[
\langle a,\pi'\rangle\deleted{\unsure{\lambda}}\ \pi'=0\qquad\text{i.e.}\quad\deleted{\unsure{\lambda}}\langle a,\pi'\rangle=0\qquad\text{i.e.}
\]
$\boxed{a\in V.}$ Conserver $f$ signifie
\[
f(a)\pi'+\psi(a)=0\tag*{(\unsure{*|8})}
\]
d'où \deleted{\unsure{nous}} en appliquant $\psi'$
\[
f(a)\underbrace{\psi'(\pi')}_{-\sigma'\pi}+\underbrace{\psi'\psi(a)}_{\beta a}=0\qquad\text{i.e.}\quad\beta a=\deleted{\unsure{h}}\sigma'f(a)\pi
\]

% archive sheet number: 5

\scanpage{7}
Si $\beta$ inversible on conclut $a=\lambda\pi$, donc l'équation qui est dessus \inserted{(\unsure{*|8})} devient
\[
\begin{aligned}
& f(a)\pi'+\lambda\underbrace{\psi(\pi)}_{-\sigma\pi'}=0\qquad\text{i.e.}\quad \sigma\lambda=\lambda^2\underbrace{f(\pi)}_{-\widetilde\sigma\widetilde\rho}\\
&\deleted{\beta\lambda\pi=\sigma'\lambda^2 f(\pi)\pi}\\
\text{i.e.}\quad &\beta\lambda=\sigma'\lambda^2\underbrace{f(\pi)}_{-\widetilde\sigma\widetilde\rho}\quad\deleted{\text{i.e.}}\qquad\text{i.e.}\quad\lambda(\sigma+\lambda\widetilde\sigma\widetilde\rho)=0\\
\text{i.e.}\quad &\lambda(\beta+\sigma'\widetilde\sigma\widetilde\rho\lambda)=0,\quad\text{en plus de }\rho\lambda=0
\end{aligned}
\]
Si on est dans un anneau intègre, celà implique \uline{soit} $\lambda=0$ i.e. $u=\mathrm{id}$, soit $\beta+\sigma'\widetilde\sigma\widetilde\rho\lambda=0$ ce qui implique (comme $\beta$ inv. par hyp) $\sigma'\widetilde\sigma\widetilde\rho$ inversible (et $\lambda=-\frac{\deleted{\unsure{\sigma'\widetilde\sigma\widetilde\rho}}\beta}{\sigma'\widetilde\sigma\widetilde\rho}$) absurde si $\rho=\bar\rho$ ($n$ impair) qui implique $\lambda=0$ (puisque $\rho\lambda=0$), sinon $\rho=2\widetilde\rho\deleted{\unsure{n}}$, si $2$ inv. \deleted{même \unsure{concl.} (\unsure{même} résultat)} même résultat. Mais si car. $2$ et $n$ pair, $\deleted{\beta\lambda=0}$ $\rho=0$ donc $\rho\lambda=0$ par des \unsure{ventres bas|ventribus|en conclure}, il reste la condition (\unsure{néc.} et suff$^{\mathrm{te}}$) $\sigma+\lambda\widetilde\sigma\widetilde\rho=0\deleted{\unsure{+}}$. Or $\sigma\sigma'=\beta+\gamma\rho=\beta$ \deleted{\unsure{donc}} inv, donc \inserted{$\sigma,\sigma'$ et} $\widetilde\sigma=\deleted{\unsure{\widetilde{\phantom\sigma}}}\sigma$ \unsure{est} inv, la condition devient $1+\lambda\widetilde\rho=0$ i.e. $\widetilde\rho$ inv. et $\lambda=-\frac1{\widetilde\rho}$. Ainsi

\uline{Prop.} Sur anneau $k$ intègre, si $\beta$ ou $\rho$ $\ne0$, alors le seul $u\in\operatorname{Aut}(E,f,\pi,\pi')$ induisant l'identité dans $V$ est l'identité, \uline{sauf} dans le cas où $\deleted{\rho=0}$, $\beta\ne0$, $n$ pair, $\deleted{\rho=2\rho'}$ car. $2$ $(\rho=2\rho'=0)$, $\rho'$ inv. auquel cas le seul autre $u$ est $\mathrm{id}\unsure{+|-}\frac1{\widetilde\rho}\pi'\otimes\pi=\mathrm{id}+\frac1{\widetilde\rho}\unsure{\langle\pi',\pi'\rangle|\langle\pi,\pi'\rangle}$ i.e. l'antipodisme

% archive sheet number: 6

\scanpage{8}
[Il reste à voir ce qui se passe si $\beta=0$, $\rho=0$ donc $\sigma\sigma'=0$. Posons
\[
\begin{aligned}
a&=\sum_0^n\lambda_i a_i,\quad\text{d'où}\quad \psi(a)=\sum-\lambda_i(\beta_0-\beta_{i-1})a_i,\\
f(a)&=\sum_0^n\lambda_i^2(\beta_0-\beta_{i-1})\unsure{-|+}\bigl(\lambda_0\lambda_1\beta_0+\lambda_1\lambda_2(\beta_0\beta_1)\deleted{+}\ldots+\lambda_{n-1}\lambda_n(\beta_0-\beta_{n-1})\bigr]
\end{aligned}
\]
\mbox{---}\qquad ]

Regardons enfin le noyau de
\[
\operatorname{Aut}(E,f,\pi,\pi')\longrightarrow\operatorname{Aut}(\overbrace{P(V)}^{C})
\]
donc les $u$ du $1^{\mathrm{er}}$ mb qui induisent une homothétie sur $V$. Cette homothétie, $\varepsilon$, conservant $f_V$, doit être dans $\mu_2$, donc $\varepsilon u$ \deleted{conserve} est l'identité sur $V$, donc de la forme
\[
\varepsilon u=\mathrm{id}_E+\deleted{\unsure{\lambda|\varepsilon}}\pi'\otimes\varepsilon a\qquad(a\in E)
\]
\deleted{et \hspace{1em}$\varepsilon\,\mathrm{id}$,}

$\deleted{u=}u=\varepsilon\,\mathrm{id}+\pi'\otimes\deleted{\varepsilon}a$ conserve $f,\pi,\pi'$.

Conserver $\pi,\pi'$
\[
\begin{aligned}
\pi&=\varepsilon\pi+\rho\deleted{\varepsilon}a\qquad\text{i.e.}\quad\boxed{(1-\varepsilon)\pi=\rho\deleted{\varepsilon}a}\\
\pi'&=\varepsilon\pi'+\deleted{\varepsilon}\langle a,\pi'\rangle\pi'\qquad\text{i.e.}\quad\deleted{\boxed{\varepsilon(1+\langle a,\pi'\rangle)=1}}\\
&&\deleted{\unsure{+a|\varepsilon a}}\boxed{\langle a,\pi'\rangle=1-\varepsilon}
\end{aligned}
\]
Enfin, conservation de $f$ \inserted{pour $u$ (ou $\varepsilon u$)} signifie
\[
\boxed{f(a)\pi'+\varepsilon\psi(a)=0}
\]

% archive sheet number: 7

\scanpage{9}
% The large diagonal strokes cancel the argument through the last scalar equation.
\deleted{Plaçons nous sur corps $k$, et supposons $\varepsilon\ne1$. (donc car. $\ne2$) -- le cas $\varepsilon=1$ étant le cas déjà traité (sous \unsure{le|la} \unsure{restriction} que $\beta,\rho$ non nuls tous deux). On doit donc avoir \uline{$\rho$ inv.} par la première équation, donc $a=\frac{\deleted{\unsure{\varepsilon|1-\varepsilon}}2}{\rho}\pi=\lambda\pi$, \deleted{donc} \inserted{$(\lambda=\frac2\rho)$} la $2^{\mathrm{ième}}$ équation donne $\rho\lambda=2\deleted{\unsure{-\varepsilon}}$ \deleted{déjà} satisfaite, enfin la dernière s'écrit}
\[
\begin{aligned}
&\deleted{\lambda^2\underbrace{f(\pi)}_{-\widetilde\rho\widetilde\sigma}\pi'\unsure{-|+}\deleted{\varepsilon}\lambda\underbrace{\psi(\pi)}_{-\sigma\pi'}=0}\\
\deleted{\text{i.e.}}\quad&\deleted{\lambda(\deleted{\unsure{-}}\underbrace{\widetilde\rho\widetilde\sigma\lambda}_{\frac{2\widetilde\rho\widetilde\sigma}{\rho}=\frac{\rho\sigma}{\rho}=\sigma}\unsure{-|+}\deleted{\varepsilon}\sigma)=0}
\end{aligned}
\]
\deleted{ce qui exige, puisque $\lambda\ne0$, que}
\[
\deleted{\widetilde\rho\widetilde\sigma\lambda\unsure{-|+}\deleted{\varepsilon}\sigma=0}
\]
\deleted{ou encore}
\[
\begin{aligned}
&\deleted{\unsure{\rho\,\tfrac{\widetilde\sigma}{2}(1-\varepsilon)}}\qquad\deleted{\sigma=\widetilde\rho\widetilde\sigma}\\
&\deleted{(\unsure{1-\varepsilon})\frac{2\widetilde\rho\widetilde\sigma}{\rho}+\varepsilon\sigma=0}
\end{aligned}
\]

Sur un anneau intègre, les seuls cas pour $\varepsilon$ sont $\varepsilon=+1$ (déjà étudié, sous la restriction que $\rho,\beta$ non tous deux inv) ou $\varepsilon=-1,\ne+\uline{1}$ (donc \deleted{car. $\ne2$)} corps des fractions $K$ de car $\ne2$) donc $\rho a=2\pi$ donc $a=\frac2\rho\pi$, $\deleted{\varepsilon}u=-\mathrm{id}+\mathord{?}\pi'\otimes\frac2\rho\pi$, c'est l'antipodisme, \inserted{\unsure{voir le Signe|vérifier le Signe}} dont on sait bien qu'il conserve $f$, et \deleted{\unsure{conserve|inverse}} $\pi,\pi'$ .. -- Donc

\uline{Proposition} \deleted{\unsure{Prop}} \inserted{(Corps de base)} Si $\beta$ ou $\rho$ est inv. \deleted{et car. $\ne2$,} alors le seul autom de $(E,f,\pi,\pi')$ qui induise l'identité

% archive sheet number: 8

\scanpage{10}
sur $C=\check{\mathbb P}(V)$, est l'\deleted{anti}identité et l'antipodisme (défini sss $\widetilde\rho$ inv \deleted{(cas $n$ pair) ou si $\rho$} comme $\underline a=\mathrm{id}_E-\frac1{\widetilde\rho}\pi'\otimes\pi$ si $n$ pair, $\underline a=\mathrm{id}_E-\frac2{\widetilde\rho=\rho}\pi'\otimes\pi$ si $n$ impair -- dans ce dernier cas, c'est d'ailleurs l'identité en car. $2$.) --

% archive sheet number: 9

\scanpage{11}
\uline{Dualité} \deleted{(\illegible{1})}

La condition $\rho$ inv est néc. et suff pour que l'accouplement naturel entre $V$ et $V'$ soit une dualité -- elle implique donc que $C$ et $C'$ sont en dualité l'un avec l'autre.

Si de plus $\beta$ inv. i.e. \deleted{\unsure{applt}} $f\mid V$ non dég. i.e. $\psi_V:V\simeq V'$, on trouve un iso entre \deleted{$V$ et} $C$ et $C'$ qui commute à $\unsure{\mathscr G|G|\mathfrak G}$ et transforme donc $Q_C$ en $Q_{C'}$. Donc dans ce cas
\[
(C,Q_C)\simeq(C',Q_{C'})
\]

Le passage d'un \deleted{\unsure{hy|hyp}}polyèdre régulier \unsure{purement|projectif} épinglé vers son dual, sous ces condition, en termes de $C^\infty_{pr\,orth}$, s'exprime donc par la conservation de $(C,Q_C)$, et en \deleted{\unsure{échangeant} \unsure{à}} \inserted{remplaçant} $\tau_0^C\ldots\tau_n^C$ par $\tau_n^C\ldots\tau_0^C$ (i.e. $\varphi_g^C$ par $\varphi_g^C\circ\theta$, où $\theta:\mathscr G\longrightarrow\mathscr G$ échange $\underline\tau_i$ et $\underline\tau_{n-i}$) -- ici il est inutile de préciser une section $a_0$ \inserted{resp $a_n$} de $C$, \deleted{\unsure{car}} puisque $\beta_0,\beta_n$ inversibles.

% archive sheet number: 10

\scanpage{12}
Soit maintenant $\Pi$ hyperpolyèdre comb. de dim $n$, $G$ le groupe de ses autom, $R_\Pi$ $G_{\deleted{\unsure{g|G}}}$-torseur des repères -- son commutant est un quotient $\unsure{\mathscr G_\Pi|G_\Pi|\mathfrak G_\Pi}$ de $\mathscr G$,

$C_{pr}(\Pi)$ cat des réa. géom régulières \inserted{proj} non dégénérées de $\Pi$ ($\Rightarrow\varphi_G:G\longrightarrow\operatorname{Aut}(X)$).
\deleted{$\unsure{C_{pr}(\Pi)}$}

A' tt $r\in R$ correspond un objet de $C_{pr}$ (cas universel épinglé) et pour $r\in R$ variable, on a un système transitif d'iso. entre ces objets. Donc $C_{pr}(\Pi)$ s'identifie \inserted{(pour $\Pi$ fixé)} à \deleted{la cat.} \deleted{des couples \unsure{$(\ )$|$(\ldots)$}} \inserted{sous cat. pleine} de $C_{pr}$ formé des $(X,\varphi_{\mathscr G},\ldots)$ tels que $\varphi_{\mathscr G}:\mathscr G\longrightarrow\operatorname{Aut}(X)$ se factorise par le quotient $\mathscr G_\Pi$ de $\mathscr G$. On fait opérer à \uline{droite} $\mathscr G_\Pi$ sur $R_\Pi$ via son action à gauche, et on l'utilise pour tordre $(X,\varphi_{\mathscr G}\ldots)$ en
\[
R_\Pi\mathbin{\wedge_{\mathscr G_\Pi}} X
\]

Même procédé pour les autres descriptions -- pour $C'_{pr}(\Pi)$ etc ....

Autre point de vue : On ne se donne pas $\Pi$, mais seulement un \inserted{groupe} quotient déterminé $\mathscr G_0$ de $\mathscr G$, \deleted{\unsure{les}} correspondant à un hyperpolyèdre régulier combinatoire épinglé. On considère un hyperpolyèdre régulier

% archive sheet number: 11

\scanpage{13}
combin. de type $\unsure{\mathscr G_0|G_0|\mathfrak G_0}$ sur $S$ comme la donnée d'un torseur \inserted{$R_0$} sous $\mathscr G_{\unsure{0S|S}}$, permettant de tordre $\Pi_0$. Une représentation de ce $\deleted{\unsure{\mathscr G}}\Pi$ revient à la donnée (en plus de $\Pi$) d'un objet de $C_{pr}$ ou d'une des cat. équivalentes. Donc la cat. des réal. géom projectives (disons) \deleted{\unsure{des|dans}} $C$ des hyperpolyèdres comb. réguliers tordus sur $S$, de type $\mathscr G_0$, est \unsure{équivalent|équivalente} à la cat. produit de $C_{pr}$ $\deleted{\unsure{(S)}}$ (ou d'un de ses avatars) et de la cat. des $\mathscr G_0$-torseurs.

\figure{\begin{tikzpicture}[x=0.65cm,y=0.65cm]
\draw[line width=.35pt] (0,0) .. controls (.25,.27) and (.33,.20) .. (.40,-.05) .. controls (.55,-.55) and (.52,-.66) .. (1.13,-.82) .. controls (1.82,-.99) and (1.93,-.92) .. (2.16,-1.35) .. controls (2.55,-2.15) and (3.34,-2.51) .. (4.04,-3.02) .. controls (4.56,-3.44) and (4.65,-3.80) .. (5.06,-4.26);
\end{tikzpicture}}

% archive sheet number: 12

\scanpage{14}
$\Pi$ $n$-\unsure{quasi|pseudo}-polyèdre comb. \uline{épinglé} universel

$G=\unsure{\mathscr G_n|\mathfrak g_n|\mathfrak G_n}$ son groupe d'autom, engendré par $(\underline\tau_i)_{0\leq i\leq n}$

$S$ schéma de base

$\boxed{C_0}={}$ catégorie \inserted{(groupoïde)} des réalisations géom non dég. projectives sur $S$
\margin{\uline{NB} $C$ catégorie rigide, dont les classes d'isom. sont en corr. 1-1 avec $\Gamma(S,\mathcal O_S)^{2n+1}$ via $\{\lambda_0,\alpha_0,\lambda_1,\alpha_1,\ldots,\lambda_{n-1},\alpha_{n-1},\lambda_n\}$}
$\simeq{}$ Catégorie des triples $(X,(\tau_*\unsure{\overset{\mathrm f}{=}|=}(\tau_i)_{0\leq i\leq n},s_0)$, $X$ fibré projectif de rg $n+1$, sur $S$, $\tau_i$ une pseudoréflexion stricte de $X$ (d'où centre $h_i$, cocentre $H_i$) $0\leq i\leq n$, $s_0$ section de $X$ sur $S$, avec conditions
\[
\left\{\begin{aligned}
1)\ &\deleted{\unsure{h_*}}(h_i)_{0\leq i\leq n}\ \deleted{\text{engendrent}}\ \inserted{\text{lin. indép.}}\ \text{en chaque fibre}\\
&(\Rightarrow C\text{ hyperplan de }X)\\
1')\ &(H_i)\ \text{lin indép en chaque fibre}\\
&(\Rightarrow c,\text{ section de }X\\
2)\ &\tau_i\tau_j=\tau_j\tau_i\quad\text{si }0\leq i,j\leq n,\ i,j\text{ non consécutifs}\\
&(\underset{\substack{\text{compte tenu}\\\text{de }1,\deleted{\unsure{1'|5}}}}{\Longleftrightarrow}\ h_i\subset H_j\ \deleted{\unsure{h_j}}\text{ si }0\leq i,j\leq n,\ i,j\text{ non conséc.})\\
3)\ &h_i\notin H_{i+1}\ \text{en tte fibre}\qquad 0\leq i\leq n-1\\
&(\Rightarrow\deleted{\unsure{\tau_i\tau_{i+1}\ne\mathrm{id}}}\ \tau_i,\tau_{i+1}\text{ ne comm. pas})\\
4)\ &s_0\subset H_i\quad(1\leq i\leq n)\\
5)\ &s_0\notin H_0\quad\text{en tte fibre}\quad\deleted{\unsure{\text{ou bien}}}
\end{aligned}\right.
\]
\margin{\uline{NB} 1') est conséquence des autres conditions 1) 2) 3) 4) 5) ; mais en l'absence d'un $s_0$ et 4) 5), il faut imposer 1') (ce qui rend 1) à 3) syst. autodual)}
% slanted
\margin{ces conditions ne faisant pas intervenir $s_0$, et \unsure{impliquent|impliquant} l'existence \unsure{locale} \unsure{unique} de $s_0$ satisfaisant 4$^\circ$)}
\margin{NB Il y a lieu d'introduire aussi $C$ cat. des couples $(E,\varphi_G^E)$ ou $(E,\tau_*)$ (\uline{sans} rig. $s_0$ ....)}

$\boxed{C''}={}$ cat des triples $(\deleted{\unsure{E}}V,\varphi_G^{\deleted{\unsure{E}}V},\deleted{\unsure{a_0}}L_0)$, $V$ fibré vectoriel de rg $n+1$ sur $S$,
\[
\varphi_G^V:G\longrightarrow\operatorname{Aut}(V)
\]
hom de grpes, $L_0\subset V$ sous-fibré inv, tels que
\[
\left\{\begin{aligned}
1)\ &\forall\,0\leq i\leq n,\ \inserted{\tau_i^V\overset{\mathrm{dfn}}{=}}\varphi_G^V(\underline\tau_i)\text{ est une ps. réflexion de }V,\\
&\text{stricte si }1\leq i\leq n,\text{ et associée à }L_0\text{ si }i=0\\
&\text{(pas néc. stricte si }i=0).\ \text{On en conclut}\\
&\deleted{\unsure{\ell}}\ \text{des }L_i\subset V,\ L'_i\subset V'\ \text{(ss-fibrés inv)}\quad 1\leq i\leq n\\
2)\ &V=\bigoplus L_i,\quad\deleted{\text{les sous-}}\bigl[\deleted{2')\ \text{des }(L'_i)_{1\leq i\leq n}\text{ indép.}}\\
&\text{en chaque fibre (i.e. }\bigcap_{1\leq i\leq n}H_i\ \text{\uline{droite}}\ \inserted{\Delta_d}\ \text{sur chaque fibre)}\bigr]\\
\deleted{3)}\ &\deleted{L_0\subset H_i\ \text{si }i\geq2}\qquad\deleted{\text{où on pose }H_i=(L'_i)^\perp}\\
&\deleted{\text{(\uline{NB} est conséquence de }1)\,2)\text{ si }\tau_0^V\text{ ps. réflex stricte)}}
\end{aligned}\right.
\]
\margin{\uline{NB} 1) 2) $\Rightarrow L_i\subset H_j$ si $1\leq i,j\leq n$ $i,j$ non cons.}
\margin{\uline{NB} 2') est conséquence des autres conditions y inclus 3), et 3) est une conséquence des autres y inclus}

% archive sheet number: 13

\scanpage{15}
$\deleted{4}\unsure{3|4})\quad L_i\notin H_{i+1}\qquad 0\leq i\leq n-1$, en tte fibre

i.e. $L_i\otimes L'_{i+1}\xrightarrow{\sim}\underline O_S$ \inserted{hom can. induit par $\deleted{\unsure{V}}V\otimes V^\vee\to\underline O_S$}

$\deleted{\unsure{C'_0}}$

\deleted{$C'_0={}$ Cat des triples précédents \uline{rigidifiés} par une section $a_0$}

\uline{NB} De 4), on tire
\[
\left\{\begin{aligned}
&L_i,L'_{i+1}\text{ en dualité}\qquad (0\leq i\leq n-1)\quad\text{donc}\\
&L_i\xrightarrow{\sim}L_{i+1}'{}^\vee\qquad (0\leq i\leq n-1)
\end{aligned}\right.
\]
De 1) on tire
\[
\left.\begin{aligned}
L'_i\otimes L_i^{\deleted{\unsure{\vee}}}&\xleftarrow{\sim}\underline O_S\\
\tau_i^V-\mathrm{id}_V&\longleftarrow\!\shortmid\,1
\end{aligned}\ \right|\quad 1\leq i\leq n\quad \deleted{\text{d'où}}\quad\text{donc}
\]
\[
L'_i\underset{\mathrm{can}}{\simeq} L_i^{\prime\vee}\qquad 1\leq i\leq n
\]
On en conclut de proche en proche des isos
\[
L_0\simeq L_1^{\prime\vee}\simeq L_1\simeq L_2^{\prime\vee}\simeq L_2\ldots\simeq L_n^{\prime\vee}\simeq L_n
\]
d'où $\backslash$
\[
V\simeq L\otimes\underline O_S^{n+1}\qquad\text{iso can.}
\]

$\boxed{C'_0}={}$ cat des triples $(V,\varphi_G^V,L_0)$ \uline{rigidifiés} par une section\inserted{-base} $a_0$ de $L_0$, i.e. des triples $(V,\varphi_G^V,a_0)$, \deleted{$a_0$ se} $V,\varphi_G^V$ comme dessus, $a_0$ section $\ne0$ en tte fibre, satisfaisant les conditions 1) à 4) où on fait $L_0=\underline O_S a_0$ [Le fait que $a_0$ est $\ne0$ sur chaque fibre est déjà \unsure{contenu} dans la première condition 4), si on l'écrit $a_0\notin H_1$ \deleted{\unsure{en}} tte fibre \ldots]

% archive sheet number: 14

\scanpage{16}
Une telle $V$ a donc une base canonique (commençant par le $a_0$ donné)
\[
(a_0,a_1,\,-\,,a_n)
\]
de telle façon que \inserted{le} $a_i$ \unsure{fasse|forme} base de $L_i$ \inserted{$0\leq i\leq n$}, \deleted{\unsure{d'où}} qu'on a
\[
\tau_i=\mathrm{id}+\deleted{\unsure{a}}a_i^\vee\otimes a_i\qquad (
\]
avec $a_i^\vee\in\Gamma V^\vee$ bien déterminé, et
\[
\langle a_i,a_{i+1}^\vee\rangle=1\qquad 0\leq i\leq n-1
\]
(On a alors exprimé déjà \deleted{1), 2)} sur $(\tau_i^V)$ 1), 2) 3), il \deleted{On aura alors \unsure{unic.|deux}} reste à exprimer \deleted{\unsure{3) et}} la commutation des $\tau_i^V,\tau_j^V$ pour $i,j$ non consécutifs, ce qui s'exprime par
\[
L_i\subset\unsure{\mathrm{Hyp}|\mathrm{Ker}}(a'_j)\qquad\text{si}\quad 0\leq i,j\leq n,\quad i,j\text{ non consécutifs.}
\]

On aura donc
\[
\left\{\begin{aligned}
 a_0^\vee&=\deleted{\unsure{(\lambda_0-1)a_0^*+\beta_0a_1^*}}(\lambda_0-1)a_0^*+\beta_0a_1^*\\
 a_i^\vee&=\deleted{\unsure{\beta|\xi}}a_{i-1}^*+(\lambda_i-1)a_i^*+\beta_i a_{i+1}^* & \unsure{0|1}\unsure{<|\leq}i\leq n-1\\
 a_n^\vee&=a_{n-1}^*+(\lambda_n-1)a_n^*
\end{aligned}\right.
\]
\margin{\uline{NB} $a_0^\vee=0$ sss $\lambda_0=1,\beta_0=0$. Donc $\tau_0$ ps. refl. stricte sss \deleted{\unsure{$\lambda_0\ne1$}} \inserted{en tt} $s\in S$, on a $\lambda_0(s)\ne1$ \uline{ou} $\beta_0(s)\ne0$, i.e. sss $(\lambda_0-1)\mathcal O_S+\beta_0\mathcal O_S=\mathcal O_S$}

\uline{Scholie} On a un foncteur $C\longrightarrow C'_0$ évident,
\[
(E,\tau_*,s_0)\longmapsto(V\deleted{\unsure{,|_{*}}}
\]
qui est une équivalence de catégories

$\boxed{C''_0}$ \deleted{\unsure{Corollaire}} Considérons la catégorie des \deleted{\unsure{couples}} \deleted{couples} triples \margin{$n\geq2$} $(C,\varphi_G^C,h_0)$, où $C$ est un fibré projectif de rg $n$, $\varphi_{\unsure{G|g}}^{C}\colon{C}\longrightarrow\mathrm{Aut}(C)$ hom de groupe, $h_0$ section de $C$, avec conditions

1) $\forall\,0\leq i\leq n$, $\tau_i^C\overset{\mathrm{dfn}}=\varphi_G(\underline{\tau}_i)$ est une pseudo-réflexion, qui est stricte si $1\leq i\leq n$, et associe un centre $h_0$ si $i=0$ (\deleted{\unsure{$\tau_0$}} pas néc. stricte) -- d'où $h_i,H_i$ $(1\leq i\leq n)$

2) Les $h_i$ $0\leq i\leq n$ lin. indép. en chaque fibre

3) $h_i\notin H_{i+1}$ \deleted{\unsure{en}} tte fibre $(0\leq i\leq n-1)$
% archive sheet number: 15

\scanpage{17}
On a un foncteur
\[
C'_0\longrightarrow C''_0
\]
C'est une équivalence de catégories \uuline{si $n\geq2$}

(mais \uline{non si $n=2$})

En effet, loc$^{\mathrm t}$ on peut écrire $C=\check{\mathbb P}(V)$, et remonter de façon \uline{unique} les \inserted{ps. réflex.} $\tau_i^C$ en des \inserted{ps. réfl.} $\tau_i^V$, qui vont satisfaire aux m$^{\mathrm s}$ conditions de commutation. \margin{\unsure{$\lambda$} (car $n\geq2$!)} Le module $V$ est déterminé modulo $\otimes L$, $L$ inv., et cette indétermination est levée par la rigidification $L_0\simeq\mathcal O_S$ \ldots

\deleted{\uline{NB}} \uline{Rescholie}

\deleted{\unsure{Donc}} Si $n\geq2$, on a \unsure{donc|dans}
\[
C_0\mathrel{\underset{\rule{1.4em}{.35pt}}{\approx}}C'_0\mathrel{\underset{\rule{1.4em}{.35pt}}{\approx}}C''_0
\]

Formes quadratiques \deleted{\unsure{pour}} invariantes sur $V$

\margin{\uuline{$n\geq2$}} \deleted{Si $n\geq$} Celles qui sont $\ne0$ sur tte fibre correspondent $1-1$ aux hyperquadriques \inserted{$Q_C$} de $C=\mathrm P(V)$, invariantes par $\varphi_{\unsure{\mathscr G|G|\mathfrak G}}^C$. On suppose ici \uline{$\lambda_i=-1$} $\forall i$ \margin{donc on remplace $\mathscr G$ par $G$, \unsure{avec|ou} les relations $\underline{\tau}_i^2=1$} % slanted
(condition néc. pour qu'il existe une $f$ \deleted{\unsure{non nulle}} inv. \unsure{avec|non|non dég.} $f_V(a_i)$ inv. \ldots) Il faut écrire
\[
\underset{\substack{\|\\c_i}}{\unsure{f_V|\varphi_V}(a_i)}\,a_i^\vee+\psi_V(a_i)=0\qquad 0\leq i\unsure{<|\leq}n
\]
i.e.
\[
\varphi_V(a_i,a_j)=-c_i\langle a_j,a_i^\vee\rangle\qquad 0\leq i,j\leq n
\]
donc
\[
\text{i)}\qquad\varphi_V(a_i,a_j)=0\qquad i,j\ \text{non consécutifs.}
\]
% archive sheet number: 16

